\documentclass[11pt,oneside]{article}

\usepackage[T1]{fontenc}
\usepackage[utf8]{inputenc}
\IfFileExists{libertinus.sty}{\usepackage{libertinus}}{\usepackage{lmodern}}
\usepackage{microtype}
\usepackage[a4paper,margin=27mm,headheight=15pt]{geometry}
\setlength{\emergencystretch}{3em}
\usepackage{amsmath,amssymb,amsthm,mathtools}
\usepackage{bm}
\usepackage{booktabs,longtable,array,tabularx}
\usepackage{enumitem}
\usepackage{graphicx}
\usepackage{xcolor}
\usepackage{listings}
\usepackage{url}
\usepackage[hidelinks]{hyperref}
\usepackage[nameinlink,noabbrev]{cleveref}
\usepackage{fancyhdr}
\usepackage{caption}
\usepackage{subcaption}

\definecolor{deepolive}{RGB}{55,67,40}
\definecolor{softkhaki}{RGB}{236,232,207}
\definecolor{softcoral}{RGB}{190,90,75}
\hypersetup{
  pdftitle={Geometric Uniform Convolutions and New Frontiers around the Fabius--Rvachev System},
  pdfsubject={A corpus-audited report on q-deformation, Thue--Morse differences, Bernoulli--Bell calculus, Fourier-zero arithmetic, sharp non-Gevrey growth, and phase-aware Lambert-W inversion},
  colorlinks=true,
  linkcolor=deepolive,
  citecolor=deepolive,
  urlcolor=softcoral
}

\pagestyle{fancy}
\fancyhf{}
\lhead{Frontiers around the Fabius--Rvachev system}
\rhead{\thepage}
\setlength{\headheight}{14pt}

% Canonical project-wide mathematical notation.
% Canonical mathematical notation for active FabiusFunction LaTeX documents.
%
% This file is intentionally a plain TeX input rather than a LaTeX package.
% Every standalone document loads it by an explicit path relative to that
% document's directory; included fragments inherit the definitions from their
% root document.  Keep this file independent of document layout and metadata.
%
% Contract:
%   * one command has one mathematical meaning throughout the active corpus;
%   * names encode semantic roles rather than a document-local choice of letter;
%   * visually similar but differently normalized objects have different print;
%   * every command below is defined normatively in the notation catalogue.
%
% The active corpus excludes docs/archive/.  Historical sources may retain
% their original notation only inside an explicitly labelled quotation.

\makeatletter
\@ifundefined{mathbb}{%
  \PackageError{fabius-notation}{The amssymb package must be loaded first}%
    {Load amsmath and amssymb before inputting fabius-notation.tex.}%
}{}
\@ifundefined{operatorname}{%
  \PackageError{fabius-notation}{The amsmath package must be loaded first}%
    {Load amsmath before inputting fabius-notation.tex.}%
}{}
\makeatother

% Version stamp used by the catalogue and by static audits.
\newcommand{\FabiusNotationVersion}{2026-08-31}

% Number systems.  NaturalNumbers is explicitly the nonnegative convention.
\newcommand{\NaturalNumbers}{\mathbb{N}_{0}}
\newcommand{\PositiveIntegers}{\mathbb{N}_{>0}}
\newcommand{\IntegerNumbers}{\mathbb{Z}}
\newcommand{\RationalNumbers}{\mathbb{Q}}
\newcommand{\RealNumbers}{\mathbb{R}}
\newcommand{\ComplexNumbers}{\mathbb{C}}
\newcommand{\ExtendedRealNumbers}{\overline{\mathbb{R}}}

% The three principal Fabius--Rvachev functions.  Subscripts are deliberate:
% a bare F and a calligraphic F have several unrelated uses in the corpus.
\newcommand{\FabiusBounded}{F_{\mathrm{bd}}}
\newcommand{\FabiusGlobal}{F_{\mathrm{gl}}}
\newcommand{\FabiusQuantile}{Q_{F}}
\newcommand{\FabiusClampedQuantile}{Q_{F,\mathrm{cl}}}
\newcommand{\RvachevUp}{\operatorname{up}}

% Constants, elementary operators, and delimiters.
\newcommand{\EulerE}{\mathrm{e}}
\newcommand{\ImaginaryUnit}{\mathrm{i}}
\newcommand{\Differential}{\mathop{}\!\mathrm{d}}
\newcommand{\AbsoluteValue}[1]{\left\lvert #1\right\rvert}
\newcommand{\Norm}[1]{\left\lVert #1\right\rVert}
\newcommand{\Floor}[1]{\left\lfloor #1\right\rfloor}
\newcommand{\Ceiling}[1]{\left\lceil #1\right\rceil}
\newcommand{\FractionalPart}[1]{\left\{#1\right\}}
\newcommand{\PositivePart}[1]{\left(#1\right)_{+}}
\newcommand{\IndicatorOf}[1]{\mathbf{1}_{#1}}
\newcommand{\SetOf}[1]{\left\{#1\right\}}
\newcommand{\InnerProduct}[1]{\left\langle #1\right\rangle}
\newcommand{\CardinalityOf}[1]{\left\lvert #1\right\rvert}
\newcommand{\DefinitionEquals}{\mathrel{\coloneqq}}
\newcommand{\Sign}{\operatorname{sgn}}
\newcommand{\IdentityOperator}{\operatorname{id}}
\newcommand{\SpanOperator}{\operatorname{span}}
\newcommand{\TraceOperator}{\operatorname{tr}}
\newcommand{\RankOperator}{\operatorname{rank}}
\newcommand{\DiagonalOperator}{\operatorname{diag}}
\newcommand{\ResidueOperator}{\operatorname*{Res}}
\newcommand{\LeastCommonMultiple}{\operatorname{lcm}}
\newcommand{\OrderOperator}{\operatorname{ord}}
\newcommand{\PrincipalLogarithm}{\operatorname{Log}}
\newcommand{\FallingFactorial}[2]{\left(#1\right)^{\underline{#2}}}
\newcommand{\RisingFactorial}[2]{\left(#1\right)^{\overline{#2}}}

% Support, topology, convolution, and metric notation.
\newcommand{\SupportOperator}{\operatorname{supp}}
\newcommand{\SupportOf}[1]{\SupportOperator\!\left(#1\right)}
\newcommand{\TopologicalSupportOf}[1]{\operatorname{tsupp}\!\left(#1\right)}
\newcommand{\Convolution}{\mathbin{*}}
\newcommand{\MetricDistance}{\operatorname{dist}}
\newcommand{\TotalVariationDistance}{d_{\mathrm{TV}}}

% Probability.  Equality and convergence in law are not metric distance.
\newcommand{\Probability}{\mathbb{P}}
\newcommand{\Expectation}{\mathbb{E}}
\newcommand{\Variance}{\operatorname{Var}}
\newcommand{\Covariance}{\operatorname{Cov}}
\newcommand{\LawOperator}{\operatorname{Law}}
\newcommand{\LawOf}[1]{\LawOperator\!\left(#1\right)}
\newcommand{\EqualInLaw}{\mathrel{\overset{\mathrm{law}}{=}}}
\newcommand{\ConvergesInLaw}{\mathrel{\xrightarrow{\mathrm{law}}}}

% Fourier and integral transforms.  The printed labels expose normalization.
% FourierTwoPi uses exp(-2*pi*i*x*xi); FourierAngular uses exp(-i*x*omega).
\newcommand{\FourierTwoPi}[1]{\widehat{#1}_{\,2\pi}}
\newcommand{\FourierAngular}[1]{\widehat{#1}_{\,\mathrm{ang}}}
\newcommand{\LaplaceTransformSymbol}{\mathcal{L}}
\newcommand{\MellinTransformSymbol}{\mathcal{M}}
\newcommand{\LaplaceTransformOf}[1]{\LaplaceTransformSymbol\!\left[#1\right]}
\newcommand{\MellinTransformOf}[1]{\MellinTransformSymbol\!\left[#1\right]}
\newcommand{\MomentGeneratingFunction}[1]{M_{#1}}
\newcommand{\CharacteristicFunction}[1]{\varphi_{#1}}

% The two standard sinc functions must never share one symbol.
\newcommand{\SincRad}{\operatorname{sinc}_{\mathrm{rad}}}
\newcommand{\SincPi}{\operatorname{sinc}_{\pi}}

% Binary arithmetic and Thue--Morse notation.
\newcommand{\BinaryDigitSum}{\operatorname{wt}_{2}}
\newcommand{\ThueMorseSignSymbol}{\varepsilon}
\newcommand{\ThueMorseSign}[1]{\varepsilon_{#1}}
\newcommand{\PadicValuation}[1]{\nu_{#1}}
\newcommand{\TwoAdicValuation}{\nu_{2}}

% q-series notation.  Every base is an explicit argument.
\newcommand{\QPochhammer}[3]{\left(#1;#2\right)_{#3}}
\newcommand{\GaussianBinomial}[3]{%
  \genfrac{[}{]}{0pt}{}{#1}{#2}_{#3}%
}
\newcommand{\QInteger}[2]{\left[#1\right]_{#2}}

% Named special functions.
\newcommand{\Polylogarithm}{\operatorname{Li}}
\newcommand{\LambertW}{W}
\newcommand{\GammaFunction}{\Gamma}
\newcommand{\RiemannZeta}{\zeta}

% Asymptotics.  BigO and LittleO are operator tokens for existing formula
% grammar; prefer the At forms so that the limiting regime is printed.
\newcommand{\BigO}{\mathcal{O}}
\newcommand{\LittleO}{o}
\newcommand{\BigOAt}[2]{\mathcal{O}_{#2}\!\left(#1\right)}
\newcommand{\LittleOAt}[2]{o_{#2}\!\left(#1\right)}

\newtheorem{theorem}{Theorem}[section]
\newtheorem{proposition}[theorem]{Proposition}
\newtheorem{lemma}[theorem]{Lemma}
\newtheorem{corollary}[theorem]{Corollary}
\newtheorem{conjecture}[theorem]{Conjecture}
\newtheorem{problem}[theorem]{Problem}
\theoremstyle{definition}
\newtheorem{definition}[theorem]{Definition}
\newtheorem{algorithm}[theorem]{Algorithm}
\theoremstyle{remark}
\newtheorem{remark}[theorem]{Remark}
\newtheorem{status}[theorem]{Status note}

\DeclareMathOperator{Var}{Var}
\DeclareMathOperator{Cov}{Cov}
\DeclareMathOperator{supp}{supp}
\DeclareMathOperator{sgn}{sgn}
\DeclareMathOperator{Li}{Li}
\DeclareMathOperator{Res}{Res}
\DeclareMathOperator{erfc}{erfc}
\newcommand{\eps}{\varepsilon}
\newcommand{\Bell}{\mathcal B}
\newcommand{\cP}{\mathcal P}
\newcommand{\cA}{\mathcal A}
\newcommand{\cB}{\mathcal B}
\newcommand{\cD}{\mathcal D}
\newcommand{\RepoCommit}{\texttt{unknown}}
\IfFileExists{repo-metadata.tex}{\input{repo-metadata.tex}}{}

\title{\textbf{Geometric Uniform Convolutions and New Frontiers\\around the Fabius--Rvachev System}\\[0.7ex]
\large A corpus-audited report on $q$-deformation, Thue--Morse differences,
Bernoulli--Bell calculus, Fourier-zero arithmetic, sharp non-Gevrey growth,
and phase-aware Lambert--$W$ inversion}
\author{Research report prepared for Vladimir Reshetnikov}
\date{30 August 2026\\Repository snapshot: \RepoCommit}

\begin{document}
\maketitle

\begin{abstract}
This report begins with a recursive full-text audit of the \texttt{*.tex} corpus in
\texttt{Analysis/FabiusFunction/docs} at repository commit \RepoCommit.  Its purpose is
not to repackage the substantial material already present there, but to develop a
coherent collection of results that are new relative to that snapshot and that connect
several of its recurring themes in a different way.

The central construction is the geometric uniform convolution
\[
 V_q=(1-q)\sum_{k\ge0}q^k Z_k,\qquad Z_k\stackrel{\mathrm{iid}}\sim\mathrm{Unif}[0,1],
 \qquad 0<q<1.
\]
Its distribution function $F_q$ and the density $u_q$ of $X_q=2V_q-1$ recover,
at $q=\tfrac12$, the Fabius function and Rvachev's up-function in their standard
probabilistic normalization.  We prove a $q$-deformed Thue--Morse formula for every
derivative of $u_q$; derive geometric Prouhet sums in closed Bell--Bernoulli form;
construct an Appell family of ``geometric Bernoulli polynomials'' that exactly inverts
averaging by $V_q$; and obtain an explicit cumulant calculus and a two-term Edgeworth
expansion in the Gaussian limit $q\uparrow1$.

For reciprocal integral bases $q=1/b$, the Fourier transform has a Hadamard product
whose zero at the integer $m$ has multiplicity $v_b(m)+1$.  This gives a
valuation-weighted zeta identity and, more strikingly, the sharp derivative-growth law
\[
 \lim_{n\to\infty}\frac{\log\|u_{1/b}^{(n)}\|_{L^p}}{n^2}
   =\frac{\log b}{2}\qquad(1\le p\le\infty).
\]
Consequently the classical up-function is smooth and compactly supported but belongs
to no Gevrey class of finite order.  Finally, a Mellin analysis of the endpoint
Laplace product exhibits a nonconstant, exponentially small log-periodic term.  Keeping
this term produces a phase-aware saddle-point formula and a stable parametric
asymptotic inversion; suppressing it recovers the familiar Lambert-$W_{-1}$ skeleton
but cannot yield a true fixed-$q$ relative asymptotic for the direct function.

Rigorous results, results conditional only on a stated saddle admissibility lemma, and
open conjectures are labeled separately.  Reproducible Python code accompanies the
report and generates all numerical tables and figures.
\end{abstract}

\tableofcontents
\clearpage

\section{Scope, corpus audit, and novelty protocol}

\subsection{What was read}

The source corpus for comparison is the complete recursive collection of files ending
in \texttt{.tex} below
\begin{center}
\path{Analysis/FabiusFunction/docs}.
\end{center}
A machine-generated manifest, including relative paths, line counts, titles when
present, and SHA-256 hashes, accompanies this report.  The scan was performed after a
fresh checkout at commit \RepoCommit.  The audit included subdirectories, drafts,
semi-formalized notes, generated articles, and auxiliary TeX fragments; no filename
filter other than the extension was applied.

The corpus already develops, often in considerable detail, the classical Fabius and
Rvachev functional equations, dyadic values, moment recurrences, infinite sinc
products, $q$-Pochhammer and $q$-binomial identities, Lambert-$W$ endpoint inversion,
Bell- and Bernoulli-polynomial machinery, interpolation on dyadic combs, and
Legendre/Lagrange expansions.  Merely deriving another form of one of those formulas
would therefore not meet the purpose of the present report.

\subsection{How overlap was controlled}

For every candidate direction we used three tests.
\begin{enumerate}[label=(\roman*)]
\item \emph{Lexical test.}  Distinctive terminology and theorem signatures were
searched across the concatenated corpus.
\item \emph{Formula test.}  Normalized equation fragments were searched after removal
of spacing, delimiter, and macro variations.
\item \emph{Structural test.}  Even when no formula matched literally, the surrounding
argument was compared against the topic and theorem inventory so that a mere change of
notation would not count as new.
\end{enumerate}
The phrase ``new'' below always means \emph{new relative to the audited repository
snapshot}.  It is not a claim of absolute priority over every item in the mathematical
literature.  The most defensible global-priority candidates are the exact quadratic
derivative-growth constant, the geometric Bernoulli inversion family, and the
phase-aware combination of the Mellin lattice with inverse-Fabius asymptotics.

\subsection{Status map}

\begin{table}[ht]
\centering
\caption{Logical status of the principal contributions.}
\begin{tabularx}{\textwidth}{@{}p{0.24\textwidth}p{0.18\textwidth}X@{}}
\toprule
Contribution & Status & Main dependency \\
\midrule
Geometric convolution $F_q,u_q$ and smoothness & proved & Fourier-product decay \\
Iterated $q$-Thue--Morse derivative formula & proved & affine-chain induction \\
Bell--Bernoulli geometric Prouhet sums & proved & exponential generating function \\
Geometric Bernoulli Appell inversion & proved & reciprocal moment-generating function \\
Cumulants and $q\uparrow1$ Edgeworth expansion & proved on compact $y$-sets & uniform characteristic-function expansion \\
Valuation-weighted Hadamard product & proved & Euler's sine product and regrouping \\
Sharp $n^2$ derivative-growth constant & proved & a resonant Fourier sequence plus refinement bounds \\
No finite Gevrey order & proved & derivative-growth contradiction \\
Mellin lattice and log-periodic Laplace asymptotic & proved & Mellin inversion and contour displacement \\
Endpoint saddle and phase-aware inversion & conditional theorem & quantified minor-arc/admissibility estimate stated in \cref{sec:saddle} \\
Legendre-envelope and transseries statements & conjectural & numerical and structural evidence \\
\bottomrule
\end{tabularx}
\end{table}

\section{A geometric-convolution deformation of the Fabius--Rvachev pair}
\label{sec:qfamily}

Throughout, $0<q<1$ and
\[
 a:=1-q,\qquad r:=\log(1/q)>0.
\]
Let $(Z_k)_{k\ge0}$ be independent random variables uniformly distributed on $[0,1]$,
and let $U_k=2Z_k-1$, so $U_k\sim\mathrm{Unif}[-1,1]$.

\begin{definition}[Geometric Fabius and up families]
Define
\begin{align}
 V_q&=a\sum_{k=0}^{\infty}q^kZ_k, & 0\le V_q\le1,\\
 X_q&=2V_q-1=a\sum_{k=0}^{\infty}q^kU_k, & -1\le X_q\le1.
\end{align}
Let
\[
 F_q(x)=\Probability(V_q\le x),\qquad f_q=F_q',
\]
and let $u_q$ denote the density of $X_q$.  Distribution functions are extended by
$F_q(x)=0$ for $x\le0$ and $F_q(x)=1$ for $x\ge1$.
\end{definition}

The almost-sure convergence is immediate from absolute convergence.  Splitting off the
first summand gives the distributional fixed points
\begin{equation}
 V_q\overset d=aZ+qV_q',\qquad X_q\overset d=aU+qX_q',
 \label{eq:fixedpoints}
\end{equation}
where primed variables are independent copies.

\begin{proposition}[Functional and symmetry equations]
The functions $F_q,f_q,u_q$ satisfy
\begin{align}
 F_q(x)+F_q(1-x)&=1,\label{eq:Fsymmetry}\\
 f_q(x)&=\frac{F_q(x/q)-F_q((x-a)/q)}{a},\label{eq:fdifference}\\
 u_q(x)&=\frac1{2a}\int_{(x-a)/q}^{(x+a)/q}u_q(y)\Differential y,\label{eq:urefinement}\\
 u_q'(x)&=\frac{u_q((x+a)/q)-u_q((x-a)/q)}{2aq}.
 \label{eq:uderivative}
\end{align}
In the left endpoint interval $0<x<a$,
\begin{equation}
 F_q'(x)=\frac1aF_q(x/q).
 \label{eq:qfabius-left}
\end{equation}
Moreover
\begin{equation}
 u_q(x)=\frac12 f_q\!\left(\frac{x+1}{2}\right).
 \label{eq:u-f-relation}
\end{equation}
At $q=a=\tfrac12$, $F_q$ is the classical Fabius function and $u_q$ is Rvachev's
up-function in the normalization $\SupportOperator u_q=[-1,1]$.
\end{proposition}

\begin{proof}
The reflection $Z_k\mapsto1-Z_k$ sends $V_q$ to $1-V_q$, proving
\eqref{eq:Fsymmetry}.  Conditioning in the first identity of \eqref{eq:fixedpoints}
and differentiating gives \eqref{eq:fdifference}; when $0<x<a$, its second term
vanishes and gives \eqref{eq:qfabius-left}.  Conditioning in the second fixed point
and changing variables gives \eqref{eq:urefinement}; differentiation yields
\eqref{eq:uderivative}.  Finally $X_q=2V_q-1$ gives \eqref{eq:u-f-relation}.  For
$q=1/2$, equations \eqref{eq:Fsymmetry} and \eqref{eq:qfabius-left} become
$F(x)+F(1-x)=1$ and $F'(x)=2F(2x)$ on $[0,1/2]$, while
\eqref{eq:uderivative} becomes
\[
 \RvachevUp'(x)=2\bigl(\RvachevUp(2x+1)-\RvachevUp(2x-1)\bigr).
\]
These equations and the endpoint normalization characterize the classical pair.
\end{proof}

\subsection{Fourier and Laplace products}

With $\SincRad z=\sin z/z$ and $\SincRad0=1$,
\begin{align}
 \CharacteristicFunction{X_q}(t)&=\Expectation e^{itX_q}
   =\prod_{k=0}^{\infty}\SincRad(aq^kt),
 \label{eq:qfourier}\\
 M_q(t)&=\Expectation e^{tV_q}
   =\prod_{k=0}^{\infty}\frac{e^{aq^kt}-1}{aq^kt},
 \label{eq:qmgf}\\
 L_q(s)&=\Expectation e^{-sV_q}
   =\prod_{k=0}^{\infty}\frac{1-e^{-aq^ks}}{aq^ks}.
 \label{eq:qlaplace}
\end{align}
Each product converges locally uniformly.  In particular, $\CharacteristicFunction{X_q}$ is entire of
exponential type at most one, as it must be for a distribution supported on
$[-1,1]$.

\begin{theorem}[Smoothness and endpoint flatness]
\label{thm:smoothness}
For every $q\in(0,1)$, the law of $X_q$ has a $C^\infty$ density $u_q$ supported on
$[-1,1]$.  The density and all its derivatives vanish at $\pm1$.  Consequently
$F_q\in C^\infty(\RealNumbers)$ after constant extension outside $[0,1]$, and every derivative
of positive order is flat at both endpoints.
\end{theorem}

\begin{proof}
Put $T=a|t|$ and $R=\log T$.  For $T\ge1$, choose
$N=\Floor{R/r}+1$.  For $0\le k<N$, $Tq^k\ge1$, and therefore
\[
 |\CharacteristicFunction{X_q}(t)|
 \le\prod_{k=0}^{N-1}\frac1{Tq^k}
 =T^{-N}q^{-N(N-1)/2}.
\]
Taking logarithms and using $N=R/r+O(1)$ gives
\begin{equation}
 \log|\CharacteristicFunction{X_q}(t)|\le-\frac{R^2}{2r}+O_q(R).
 \label{eq:lognormal-fourier-bound}
\end{equation}
Hence $|t|^m\CharacteristicFunction{X_q}(t)\in L^1(\RealNumbers)$ for every $m$.  Fourier inversion gives a
$C^\infty$ density and permits termwise differentiation to all orders.  Since the
random series lies in $[-1,1]$, this density has compact support; smooth extension by
zero forces all endpoint jets to vanish.  Equation \eqref{eq:u-f-relation} transfers
the assertion to $F_q$.
\end{proof}

\begin{remark}
The decay in \eqref{eq:lognormal-fourier-bound} is faster than every inverse power but
much slower than $e^{-c|t|^\alpha}$.  This already hints that the functions are smooth
but far outside familiar finite-order Gevrey scales.  In \cref{sec:gevrey} we prove a
sharp statement, not merely a heuristic one, for $q=1/b$.
\end{remark}

\begin{proposition}[Log-concavity]
For each $q\in(0,1)$, $u_q$ is even and log-concave on $(-1,1)$.  Thus it is unimodal
with its maximum at zero; $f_q$ is symmetric about $1/2$ and log-concave on $(0,1)$.
\end{proposition}

\begin{proof}
Every finite partial sum of $X_q$ has a density obtained by convolving scaled uniform
densities.  Uniform densities on intervals are log-concave, and convolution preserves
log-concavity.  The finite densities converge weakly to $u_q$ and, by
\cref{thm:smoothness}, locally uniformly after passing through Fourier inversion.
Log-concavity is closed under such a nonzero limit.  Evenness follows from
$U_k\overset d=-U_k$.  The statement for $f_q$ follows by affine rescaling.
\end{proof}

\section{A $q$-Thue--Morse calculus for all derivatives}
\label{sec:qthue}

The two affine branches in \eqref{eq:uderivative} generate a binary tree.  Their signs
are precisely the signs that become the Thue--Morse sequence when $q=1/2$, while their
locations form a geometric rather than an arithmetic Prouhet configuration.

For a sign word $\boldsymbol\epsilon=(\epsilon_0,\dots,\epsilon_{n-1})\in\{\pm1\}^n$,
write
\[
 \epsilon(\boldsymbol\epsilon)=\prod_{j=0}^{n-1}\epsilon_j,
 \qquad
 Q_q(\boldsymbol\epsilon)=\sum_{j=0}^{n-1}\epsilon_jq^j.
\]

\begin{theorem}[Iterated geometric Thue--Morse formula]
\label{thm:qthue-derivative}
For every integer $n\ge1$,
\begin{equation}
 u_q^{(n)}(x)=
 \frac{q^{-n(n-1)/2}}{(2aq)^n}
 \sum_{\boldsymbol\epsilon\in\{\pm1\}^n}
 \epsilon(\boldsymbol\epsilon)\,
 u_q\!\left(q^{-n}\bigl(x+aQ_q(\boldsymbol\epsilon)\bigr)\right).
 \label{eq:qthue-derivative}
\end{equation}
At $q=1/2$, if $\ThueMorseSignSymbol_j=(-1)^{\BinaryDigitSum(j)}$ is the $\{\pm1\}$ Thue--Morse sequence, this
reduces to
\begin{equation}
 \RvachevUp^{(n)}(x)=(-1)^n2^{n(n+1)/2}
 \sum_{j=0}^{2^n-1}\ThueMorseSignSymbol_j\,
 \RvachevUp\bigl(2^nx+2j-(2^n-1)\bigr).
 \label{eq:classical-thue-derivative}
\end{equation}
\end{theorem}

\begin{proof}
The case $n=1$ is \eqref{eq:uderivative}.  Suppose \eqref{eq:qthue-derivative} holds.
Differentiating a term introduces the slope $q^{-n}$, and substituting
\eqref{eq:uderivative} for the resulting $u_q'$ appends one sign and one affine branch.
Thus the scalar multiplier acquires $(2aq)^{-1}q^{-n}$, changing
$q^{-n(n-1)/2}$ to $q^{-(n+1)n/2}$.  Reversal of the order of signs does not change
either their product or the set of geometric sums, completing the induction.

For $q=1/2$, the numbers
$2^{n-1}Q_{1/2}(\boldsymbol\epsilon)$ are all odd integers from
$-(2^n-1)$ to $2^n-1$.  Writing $\epsilon_j=2\delta_j-1$ shows that the product of
signs is $(-1)^n(-1)^{\sum\delta_j}$; bit reversal preserves the binary digit sum.
This gives \eqref{eq:classical-thue-derivative}.
\end{proof}

\subsection{Geometric Prouhet sums}

Define
\begin{equation}
 S_{n,m}(q)=\sum_{\boldsymbol\epsilon\in\{\pm1\}^n}
 \epsilon(\boldsymbol\epsilon)Q_q(\boldsymbol\epsilon)^m.
 \label{eq:Snm-def}
\end{equation}
These are the signed moments of the branch configuration in
\eqref{eq:qthue-derivative}.

\begin{theorem}[Bell--Bernoulli evaluation]
\label{thm:geometric-prouhet}
The exponential generating function is
\begin{equation}
 \sum_{m\ge0}S_{n,m}(q)\frac{t^m}{m!}
 =2^n\prod_{j=0}^{n-1}\sinh(tq^j).
 \label{eq:Snm-egf}
\end{equation}
Consequently $S_{n,m}(q)=0$ when $m<n$ or $m\not\equiv n\pmod2$.  For $m=n+2h$,
put
\begin{equation}
 \alpha_\ell(n,q)=
 \frac{2^{2\ell-1}B_{2\ell}}{\ell(2\ell)!}
 \frac{1-q^{2\ell n}}{1-q^{2\ell}}
 \qquad(\ell\ge1).
 \label{eq:alpha-def}
\end{equation}
If $Y_h$ denotes the complete exponential Bell polynomial, then
\begin{equation}
 S_{n,n+2h}(q)=
 2^n(n+2h)!q^{n(n-1)/2}
 \frac{Y_h(1!\alpha_1,2!\alpha_2,\dots,h!\alpha_h)}{h!}.
 \label{eq:Snm-bell}
\end{equation}
In particular,
\begin{align}
 S_{n,n}(q)&=2^n n!q^{n(n-1)/2},\\
 S_{n,n+2}(q)&=
 \frac{2^n(n+2)!}{6}q^{n(n-1)/2}
 \frac{1-q^{2n}}{1-q^2}.
\end{align}
\end{theorem}

\begin{proof}
Summing one sign at a time gives
\[
 \sum_{\boldsymbol\epsilon}\epsilon(\boldsymbol\epsilon)
 e^{tQ_q(\boldsymbol\epsilon)}
 =\prod_{j=0}^{n-1}(e^{tq^j}-e^{-tq^j}),
\]
which is \eqref{eq:Snm-egf}.  Factor $tq^j$ out of every hyperbolic sine and use
\begin{equation}
 \log\frac{\sinh z}{z}
 =\sum_{\ell\ge1}
 \frac{2^{2\ell-1}B_{2\ell}}{\ell(2\ell)!}z^{2\ell}.
 \label{eq:logsinh-bernoulli}
\end{equation}
Then
\[
 2^n\prod_{j=0}^{n-1}\sinh(tq^j)
 =2^nt^nq^{n(n-1)/2}
 \exp\!\left(\sum_{\ell\ge1}\alpha_\ell(n,q)t^{2\ell}\right).
\]
The defining generating function of the complete Bell polynomials yields
\eqref{eq:Snm-bell}.
\end{proof}

\begin{corollary}[Explicit centered Thue--Morse power sums]
For the ordinary Thue--Morse signs $\ThueMorseSignSymbol_j=(-1)^{\BinaryDigitSum(j)}$,
\begin{equation}
 \sum_{j=0}^{2^n-1}\ThueMorseSignSymbol_j\bigl(2j-(2^n-1)\bigr)^m
 =(-1)^n2^{(n-1)m}S_{n,m}(1/2).
 \label{eq:centered-TM-sum}
\end{equation}
Thus \eqref{eq:Snm-bell} gives a closed formula for every centered power sum, not only
the classical vanishing range $m<n$.
\end{corollary}

\begin{remark}
Equation \eqref{eq:Snm-bell} is simultaneously a $q$-deformation of Prouhet
cancellation and a finite-level shadow of the Bell--Bernoulli cumulant formulas below.
The same coefficient
$(1-q^{2\ell n})/(1-q^{2\ell})$ is the finite geometric power sum that becomes
$(1-q^{2\ell})^{-1}$ in the infinite convolution.
\end{remark}

\section{Cumulants, moments, and geometric Bernoulli polynomials}
\label{sec:bernoulli}

\subsection{Exact cumulants}

\begin{theorem}[Bernoulli cumulants]
\label{thm:cumulants}
The nonzero cumulants of $X_q$ are
\begin{equation}
 \kappa_{2m}(X_q)=
 \frac{2^{2m}B_{2m}}{2m}\,
 \frac{a^{2m}}{1-q^{2m}},
 \qquad m\ge1,
 \label{eq:X-cumulants}
\end{equation}
and all odd cumulants vanish.  For $V_q$,
\begin{equation}
 \kappa_1(V_q)=\frac12,\qquad
 \kappa_{2m}(V_q)=
 \frac{B_{2m}}{2m}\,
 \frac{a^{2m}}{1-q^{2m}},
 \qquad
 \kappa_{2m+1}(V_q)=0\quad(m\ge1).
 \label{eq:V-cumulants}
\end{equation}
Consequently every raw or centered moment is an explicit complete Bell polynomial in
rational functions of $q$ and Bernoulli numbers.
\end{theorem}

\begin{proof}
For $U\sim\mathrm{Unif}[-1,1]$,
\[
 \log\Expectation e^{tU}=\log\frac{\sinh t}{t},
\]
and \eqref{eq:logsinh-bernoulli} shows that
$\kappa_{2m}(U)=2^{2m}B_{2m}/(2m)$.  Cumulants add under independence and scale
homogeneously, so summing $(aq^k)^{2m}$ proves \eqref{eq:X-cumulants}.
Since $V_q=(X_q+1)/2$, \eqref{eq:V-cumulants} follows.  The moment--cumulant relation
is precisely the complete Bell-polynomial identity.
\end{proof}

For example,
\begin{align}
 \Variance(V_q)&=\frac{a}{12(1+q)},\label{eq:Vvariance}\\
 \kappa_4(V_q)&=-\frac{a^4}{120(1-q^4)},\\
 \kappa_6(V_q)&=\frac{a^6}{252(1-q^6)}.
\end{align}
The moments also obey a triangular recurrence obtained directly from
\eqref{eq:fixedpoints}.  If $m_n(q)=\Expectation V_q^n$, then $m_0=1$ and
\begin{equation}
 (1-q^n)m_n(q)=
 \sum_{j=0}^{n-1}\binom nj
 \frac{q^j a^{n-j}}{n-j+1}m_j(q).
 \label{eq:moment-recurrence}
\end{equation}
For rational $q$, all moments are rational.

\subsection{An Appell family that inverts geometric averaging}

Classical Bernoulli polynomials invert averaging over one uniform random variable:
$\Expectation[B_n(x+Z)]=x^n$.  The entire geometric convolution has a natural analogue.

\begin{definition}[Geometric Bernoulli polynomials]
Define $\Bell_{n,q}(x)$ by
\begin{equation}
 \frac{e^{xt}}{M_q(t)}
 =\prod_{k=0}^{\infty}\frac{aq^kt}{e^{aq^kt}-1}\,e^{xt}
 =\sum_{n=0}^{\infty}\Bell_{n,q}(x)\frac{t^n}{n!}.
 \label{eq:geometric-Bernoulli-egf}
\end{equation}
The product converges near $t=0$, and the definition may equally be read as a formal
power-series identity.
\end{definition}

\begin{theorem}[Averaging inverse and $q$-difference law]
\label{thm:geometric-Bernoulli}
The polynomials $\Bell_{n,q}$ form an Appell sequence and satisfy
\begin{align}
 \frac{\Differential}{\Differential x}\Bell_{n,q}(x)&=n\Bell_{n-1,q}(x),
 \label{eq:appell}\tag{A}\\
 \Expectation\,\Bell_{n,q}(x+V_q)&=x^n,
 \label{eq:averaging-inverse}\tag{B}\\
 \Bell_{n,q}(x+a)-\Bell_{n,q}(x)
 &=anq^{n-1}\Bell_{n-1,q}(x/q),
 \label{eq:q-difference-Bernoulli}\tag{C}\\
 \Bell_{n,q}(1-x)&=(-1)^n\Bell_{n,q}(x).
 \label{eq:Bernoulli-reflection}\tag{D}
\end{align}
As $q\downarrow0$, $\Bell_{n,q}(x)$ tends coefficientwise to the classical Bernoulli
polynomial $B_n(x)$.
\end{theorem}

\begin{proof}
Equation \eqref{eq:appell} follows by differentiating
\eqref{eq:geometric-Bernoulli-egf}.  Multiplying its generating function by
$M_q(t)=\Expectation e^{tV_q}$ proves \eqref{eq:averaging-inverse}.  The fixed-point product
satisfies
\[
 M_q(t)=\frac{e^{at}-1}{at}M_q(qt).
\]
Therefore, with $G_q(x,t)=e^{xt}/M_q(t)$,
\[
 G_q(x+a,t)-G_q(x,t)=at\,G_q(x/q,qt),
\]
and coefficient extraction gives \eqref{eq:q-difference-Bernoulli}.  Symmetry
$V_q\overset d=1-V_q$ is equivalent to $M_q(t)=e^tM_q(-t)$; hence
$G_q(1-x,t)=G_q(x,-t)$, proving \eqref{eq:Bernoulli-reflection}.  At $q=0$, only the
first uniform summand remains and the generating function becomes
$te^{xt}/(e^t-1)$.
\end{proof}

Writing $y=x-1/2$, $\sigma_q^2=\Variance(V_q)$, and $\kappa_{4,q}=\kappa_4(V_q)$ gives
\begin{align}
 \Bell_{0,q}(x)&=1,\\
 \Bell_{1,q}(x)&=y,\\
 \Bell_{2,q}(x)&=y^2-\sigma_q^2,\\
 \Bell_{3,q}(x)&=y^3-3\sigma_q^2y,\\
 \Bell_{4,q}(x)&=y^4-6\sigma_q^2y^2+3\sigma_q^4-\kappa_{4,q}.
\end{align}
More generally, $\Bell_{n,q}(x)$ is the complete Bell polynomial whose first argument
is $y$ and whose higher arguments are the negatives of the centered cumulants of
$V_q$.

\begin{corollary}[Exact moment inversion]
For $m_j(q)=\Expectation V_q^j$,
\begin{equation}
 \sum_{j=0}^{n}\binom njm_j(q)\Bell_{n-j,q}(x)=x^n.
 \label{eq:moment-inversion}
\end{equation}
Thus the lower-triangular moment matrix of the geometric convolution has an explicit
Appell-polynomial inverse.
\end{corollary}

\begin{remark}
At finite depth $N$, the same construction is a multiple Bernoulli/N\o rlund
polynomial with geometric step sizes $a,aq,\dots,aq^{N-1}$.  The infinite-depth limit
retains an exact one-step $q$-difference law, which is not present for arbitrary
multiple Bernoulli systems.  This provides a direct bridge from Bernoulli polynomials
to the probabilistic Fabius fixed point rather than passing only through moment
coefficients.
\end{remark}

\section{The Gaussian end of the deformation: a mod-Gaussian expansion}
\label{sec:gaussian}

As $q\uparrow1$, the support remains fixed but the distribution concentrates at its
center because the individual geometric weights become small.  Normalize by
\begin{equation}
 Y_q=\frac{V_q-1/2}{\sigma_q},
 \qquad
 \sigma_q^2=\frac{1-q}{12(1+q)}.
 \label{eq:Yq-def}
\end{equation}
Let $\lambda_j(q)$ denote the standardized cumulants of $Y_q$.

\begin{proposition}[Standardized cumulants]
For $m\ge2$,
\begin{equation}
 \lambda_{2m}(q)=
 \frac{B_{2m}}{2m}\frac{a^{2m}}{1-q^{2m}}
 \left(\frac{12(1+q)}a\right)^m
 =O_m(a^{m-1}),
 \label{eq:standardized-cumulants}
\end{equation}
and the odd cumulants vanish.  In particular,
\begin{align}
 \lambda_4(q)&=-\frac65\frac{1-q^2}{1+q^2}
 =-\frac65a-\frac35a^2+O(a^3),\label{eq:lambda4}\\
 \lambda_6(q)&=\frac{48}{7}\frac{a^2(1+q)^3}{1+q+q^2+q^3+q^4+q^5}
 =\frac{64}{7}a^2+O(a^3).
 \label{eq:lambda6}
\end{align}
\end{proposition}

\begin{theorem}[Gaussian and second-order Edgeworth limits]
\label{thm:edgeworth}
Let $g_q$ and $G_q$ be the density and distribution function of $Y_q$, and let
$\phi,\Phi$ be the standard normal density and distribution function.  For $y$ in a
fixed compact set, as $a=1-q\downarrow0$,
\begin{align}
 g_q(y)=\phi(y)\bigg[1
 &-\frac a{20}H_4(y)\\
 &+a^2\left(-\frac1{40}H_4(y)+\frac4{315}H_6(y)
                 +\frac1{800}H_8(y)\right)
 +O(a^3)\bigg],
 \label{eq:edgeworth-density}
\end{align}
where $H_n$ are the probabilists' Hermite polynomials.  Integrating,
\begin{align}
 G_q(y)=\Phi(y)+\phi(y)\bigg[&\frac a{20}H_3(y)\\
 &+a^2\left(\frac1{40}H_3(y)-\frac4{315}H_5(y)
                   -\frac1{800}H_7(y)\right)
 +O(a^3)\bigg].
 \label{eq:edgeworth-cdf}
\end{align}
In particular, $Y_q\Rightarrow N(0,1)$.
\end{theorem}

\begin{proof}
The logarithm of the characteristic function is
\[
 \log\Expectation e^{itY_q}=-\frac{t^2}{2}
 +\sum_{m\ge2}\lambda_{2m}(q)\frac{(it)^{2m}}{(2m)!}.
\]
Use \eqref{eq:lambda4}--\eqref{eq:lambda6}, exponentiate through order $a^2$, and
Fourier-invert.  The $\lambda_4^2/2$ term produces $H_8/800$.  The product structure
and the bound used in \cref{thm:smoothness}, after standardization, provide an
integrable majorant outside a growing neighborhood of the origin; this makes the
expansion uniform on compact $y$-sets.  Finally
$\int_{-\infty}^y\phi(t)H_n(t)\Differential t=-\phi(y)H_{n-1}(y)$ gives
\eqref{eq:edgeworth-cdf}.
\end{proof}

\begin{remark}
The limit is a triangular-array central limit theorem with a particularly transparent
small parameter: the largest normalized summand is $O(\sqrt{1-q})$.  Formula
\eqref{eq:standardized-cumulants} supplies an all-orders mod-Gaussian expansion.  At
order $a^N$, only finitely many Bernoulli cumulants contribute, and the coefficient is
an explicit complete Bell polynomial in Hermite operators.
\end{remark}

\IfFileExists{figures/edgeworth-density.pdf}{
\begin{figure}[ht]
\centering
\includegraphics[width=0.78\textwidth]{figures/edgeworth-density.pdf}
\caption{A numerical comparison generated by \texttt{experiments.py}: standardized
$q$-convolution density, Gaussian approximation, and the first two Edgeworth
corrections.}
\label{fig:edgeworth}
\end{figure}}{}

\section{Fourier zeros, valuations, and weighted zeta identities}
\label{sec:zeros}

It is useful first to remove the harmless support normalization.  Define
\begin{equation}
 \Psi_q(z)=\prod_{k=0}^{\infty}\SincRad(q^kz).
 \label{eq:Psiq}
\end{equation}
Then $\CharacteristicFunction{X_q}(t)=\Psi_q(at)$.  The zero multiset is
\begin{equation}
 \mathcal Z_q=\{\,\pi n q^{-k}:n\in\IntegerNumbers\setminus\{0\},\ k\ge0\,\},
 \label{eq:zero-multiset}
\end{equation}
with multiplicity equal to the number of representations.

\begin{proposition}[Arithmetic dichotomy for zero collisions]
All nonzero zeros of $\Psi_q$ are simple if and only if
$q^d\notin\RationalNumbers$ for every integer $d\ge1$.  If $q^d\in\RationalNumbers$ for some
$d\ge1$, then $\Psi_q$ has infinitely many multiple zeros.
\end{proposition}

\begin{proof}
Two factors share a zero exactly when
$nq^{-k}=mq^{-\ell}$ for nonzero integers $m,n$ and $k\ne\ell$.  This is equivalent to
$q^{|k-\ell|}\in\RationalNumbers$.  One collision can be multiplied by arbitrary nonzero
integers, yielding infinitely many.
\end{proof}

The cleanest arithmetic occurs at reciprocal integral bases.

\begin{theorem}[Valuation-weighted Hadamard product]
\label{thm:hadamard-valuation}
Let $b\ge2$ be an integer and let
\[
 v_b(m)=\max\{j\ge0:b^j\mid m\}.
\]
Then
\begin{equation}
 \Psi_{1/b}(z)=
 \prod_{m=1}^{\infty}
 \left(1-\frac{z^2}{\pi^2m^2}\right)^{v_b(m)+1}.
 \label{eq:valuation-hadamard}
\end{equation}
The product converges locally uniformly.  Equivalently, the zero at $\pm\pi m$ has
multiplicity $v_b(m)+1$ and there are no other zeros.
\end{theorem}

\begin{proof}
Euler's product gives
\[
 \SincRad(z/b^k)=\prod_{n=1}^{\infty}
 \left(1-\frac{z^2}{\pi^2n^2b^{2k}}\right).
\]
A factor indexed by $(n,k)$ contributes to the zero $\pi m$ exactly when
$m=nb^k$, and the number of such $k$ is $v_b(m)+1$.  Regrouping is legitimate because
on compact sets the sum of absolute quadratic deviations is bounded by
\[
 C|z|^2\sum_{k\ge0}b^{-2k}\sum_{n\ge1}n^{-2}<\infty.
\]
\end{proof}

\begin{corollary}[Valuation Dirichlet series]
For $\Re s>1$,
\begin{equation}
 \sum_{m=1}^{\infty}\frac{v_b(m)+1}{m^s}
 =\frac{\zeta(s)}{1-b^{-s}}.
 \label{eq:valuation-zeta}
\end{equation}
Consequently, for $|z|<\pi$,
\begin{equation}
 \log\Psi_{1/b}(z)
 =-\sum_{j=1}^{\infty}
 \frac{\zeta(2j)}{j(1-b^{-2j})}
 \left(\frac z\pi\right)^{2j}.
 \label{eq:Psi-Taylor-zeta}
\end{equation}
\end{corollary}

\begin{proof}
Since $v_b(m)+1=\sum_{k\ge0}\IndicatorOf{b^k\mid m}$,
\[
 \sum_{m\ge1}\frac{v_b(m)+1}{m^s}
 =\sum_{k\ge0}\sum_{n\ge1}(b^kn)^{-s}.
\]
This proves \eqref{eq:valuation-zeta}; logarithmic expansion of
\eqref{eq:valuation-hadamard} gives \eqref{eq:Psi-Taylor-zeta}.
\end{proof}

\begin{corollary}[A cotangent/valuation partial fraction identity]
Away from the zeros,
\begin{equation}
 \sum_{k=0}^{\infty}
 \left(b^{-k}\cot(z/b^k)-\frac1z\right)
 =-2z\sum_{m=1}^{\infty}
 \frac{v_b(m)+1}{\pi^2m^2-z^2}.
 \label{eq:cotangent-valuation}
\end{equation}
Both sides converge normally on compact subsets avoiding the poles.
\end{corollary}

\begin{remark}
For $b=2$, the multiplicity is $v_2(m)+1$, so the zero divisor of the up-function's
Fourier transform literally records the dyadic valuation.  Equation
\eqref{eq:valuation-zeta} says that the corresponding spectral zeta function is the
ordinary Riemann zeta function multiplied by the local Euler correction
$(1-2^{-s})^{-1}$.  This is a direct arithmetic encoding in the sinc product, distinct
from the more familiar appearance of Thue--Morse signs in real-space derivatives.
\end{remark}

\section{Sharp derivative growth and failure of every finite Gevrey bound}
\label{sec:gevrey}

The superpolynomial Fourier decay in \eqref{eq:lognormal-fourier-bound} proves
smoothness but does not identify the exact derivative scale.  Reciprocal integral bases
admit a resonant sequence on which the Fourier product can be evaluated sharply enough
to close the upper and lower bounds.

Let
\begin{equation}
 S_b=\sum_{k=0}^{\infty}b^{-k}U_k
 \label{eq:Sb}
\end{equation}
with density $h_b$.  Thus $X_{1/b}=\frac{b-1}{b}S_b$, so $h_b$ and $u_{1/b}$ differ
only by a fixed dilation.

\begin{lemma}[Refinement upper bound]
\label{lem:derivative-upper}
For every $n\ge0$,
\begin{equation}
 \|h_b^{(n)}\|_\infty
 \le b^{n(n+1)/2}\|h_b\|_\infty.
 \label{eq:derivative-upper}
\end{equation}
\end{lemma}

\begin{proof}
Since $S_b\overset d=U+b^{-1}S_b'$, conditioning on $U$ gives
\[
 h_b(x)=\frac12\int_{b(x-1)}^{b(x+1)}h_b(y)\Differential y,
\]
so
\[
 h_b'(x)=\frac b2\bigl(h_b(b(x+1))-h_b(b(x-1))\bigr).
\]
After $n$ iterations the common scalar before the signed $2^n$-term sum is
$(b/2)^n b^{n(n-1)/2}$.  Taking absolute values cancels the $2^{-n}$ and yields
\eqref{eq:derivative-upper}.
\end{proof}

\begin{lemma}[Exact resonant Fourier values]
\label{lem:resonant-values}
Put
\[
 c_b=\frac{\pi}{b+1},\qquad
 \theta_b=\frac{\sin c_b}{c_b},\qquad
 T_b=\prod_{j=1}^{\infty}\SincRad(c_b/b^j)>0.
\]
Then, for every $m\ge0$,
\begin{equation}
 \left|\CharacteristicFunction{S_b}(c_bb^m)\right|
 =T_b\theta_b^{m+1}b^{-m(m+1)/2}.
 \label{eq:resonant-fourier}
\end{equation}
\end{lemma}

\begin{proof}
We have $\CharacteristicFunction{S_b}=\Psi_{1/b}$.  For $0\le j\le m$, the congruence
$b^j\equiv(-1)^j\pmod{b+1}$ gives
$|\sin(c_bb^j)|=\sin c_b$.  Hence
\[
 \prod_{j=0}^{m}|\SincRad(c_bb^j)|
 =\theta_b^{m+1}b^{-m(m+1)/2}.
\]
The remaining factors are exactly $T_b$ after reindexing.
\end{proof}

\begin{theorem}[Sharp quadratic derivative growth]
\label{thm:sharp-derivative-growth}
For every integer $b\ge2$ and every $1\le p\le\infty$,
\begin{equation}
 \boxed{
 \lim_{n\to\infty}
 \frac{\log\|u_{1/b}^{(n)}\|_{L^p(\RealNumbers)}}{n^2}
 =\frac{\log b}{2}.}
 \label{eq:sharp-derivative-growth}
\end{equation}
Equivalently,
\[
 \|u_{1/b}^{(n)}\|_p=b^{n^2/2+O_{b,p}(n)}.
\]
For the classical up-function,
\begin{equation}
 \log_2\|\RvachevUp^{(n)}\|_p=\frac{n^2}{2}+O_p(n).
 \label{eq:up-derivative-growth}
\end{equation}
\end{theorem}

\begin{proof}
The upper bound for $p=\infty$ follows from \cref{lem:derivative-upper}; compact
support then gives the same logarithmic upper bound for every $p$.  For the lower
bound, Fourier transformation gives
\[
 |t|^n|\CharacteristicFunction{S_b}(t)|=|\FourierAngular{h_b^{(n)}}(-t)|
 \le\|h_b^{(n)}\|_1.
\]
Choose $t=c_bb^n$ in \eqref{eq:resonant-fourier}.  Its logarithm is
\[
 n\log(c_bb^n)+(n+1)\log\theta_b
 -\frac{n(n+1)}2\log b+\log T_b
 =\frac{\log b}{2}n^2+O_b(n).
\]
Thus the $L^1$ norm has the matching lower bound.  On a fixed compact support,
$\|f\|_1\le C_p\|f\|_p$ for all $p\ge1$, so the lower bound transfers to every
$L^p$ norm.  Finally, the fixed dilation between $h_b$ and $u_{1/b}$ changes the
logarithm of an $n$th derivative norm by only $O_b(n)$.
\end{proof}

\begin{corollary}[No finite Gevrey order]
\label{cor:no-gevrey}
For $b\ge2$, neither $u_{1/b}$ nor the associated $F_{1/b}$ belongs to a Gevrey class
$G^s$ of any finite order $s$ on a neighborhood of its closed support.  In particular,
the classical Fabius and up-functions are not Gevrey of finite order.
\end{corollary}

\begin{proof}
A compactly supported $G^s$ function satisfies
$\|f^{(n)}\|_\infty\le CR^n(n!)^s$, whose logarithm is $O_s(n\log n)$.  This
contradicts \eqref{eq:sharp-derivative-growth}.  The affine and differential relation
between $F_{1/b}$ and $u_{1/b}$ transfers the conclusion.
\end{proof}

\begin{remark}
The result is sharper than the qualitative statement ``nowhere analytic.''  It locates
the precise Denjoy--Carleman scale: the natural derivative majorant is
$M_n\asymp b^{n^2/2+O(n)}$, not $(n!)^s$ for any fixed $s$.  Since
$M_{n-1}/M_n$ decays geometrically, the associated class is strongly nonquasianalytic,
consistent with compact support.
\end{remark}

\IfFileExists{figures/resonant-fourier.pdf}{
\begin{figure}[ht]
\centering
\includegraphics[width=0.78\textwidth]{figures/resonant-fourier.pdf}
\caption{Exact resonant Fourier values from \eqref{eq:resonant-fourier}, displayed on
the quadratic logarithmic scale that drives \cref{thm:sharp-derivative-growth}.}
\label{fig:resonant}
\end{figure}}{}

\section{Mellin analysis and the hidden log-periodic phase}
\label{sec:mellin}

The left endpoint of $F_q$ is controlled by the Laplace transform
\eqref{eq:qlaplace}.  Set
\begin{equation}
 H_q(z)=\prod_{k=0}^{\infty}\frac{1-e^{-zq^k}}{zq^k},
 \qquad L_q(s)=H_q(as),
 \label{eq:Hq}
\end{equation}
and write $L=\log z$.  The discrete scale $z\mapsto qz$ produces a Mellin lattice on
the imaginary axis.  Its contribution is nonzero but, for $q=1/2$, exponentially
small enough to be easily missed numerically.

Let $\gamma$ be Euler's constant and $\gamma_1$ the first Stieltjes constant in the
convention
\[
 \zeta(1+s)=\frac1s+\sum_{n\ge0}\frac{(-1)^n\gamma_n}{n!}s^n.
\]
Define
\begin{equation}
 c_1=\frac{\pi^2}{12}-\frac{\gamma^2}{2}-\gamma_1,
 \qquad
 \chi_n=\frac{2\pi in}{r},
 \label{eq:c1-chi}
\end{equation}
and the real-analytic $r$-periodic function
\begin{equation}
 \cP_r(L)=
 -\frac1r\sum_{n\in\IntegerNumbers\setminus\{0\}}
 \frac{\Gamma(1+\chi_n)\zeta(1+\chi_n)}{\chi_n}
 e^{-\chi_nL}.
 \label{eq:periodic-P}
\end{equation}
The coefficients occur in conjugate pairs.  Stirling's formula gives exponential
convergence, with the $n$th harmonic of order
$e^{-\pi^2|n|/r}$ up to powers of $|n|$.

\begin{lemma}[Mellin kernel]
For
\[
 g(x)=\log\frac{1-e^{-x}}x,
\]
its Mellin transform in $-1<\Re s<0$ is
\begin{equation}
 g^*(s)=\int_0^\infty g(x)x^{s-1}\Differential x
 =-\frac{\Gamma(1+s)\zeta(1+s)}s.
 \label{eq:mellin-kernel}
\end{equation}
Moreover,
\begin{equation}
 \log H_q(z)=\frac1{2\pi i}\int_{c-i\infty}^{c+i\infty}
 g^*(s)\frac{z^{-s}}{1-e^{rs}}\Differential s,
 \qquad -1<c<0.
 \label{eq:mellin-H}
\end{equation}
\end{lemma}

\begin{proof}
Differentiate $g$:
$g'(x)=(e^x-1)^{-1}-x^{-1}$.  Integration by parts is legitimate in the stated
strip and analytic continuation of the classical Bose integral gives
\[
 \int_0^\infty g'(x)x^s\Differential x=\Gamma(1+s)\zeta(1+s).
\]
This proves \eqref{eq:mellin-kernel}.  Mellin inversion applied to
$\sum_{k\ge0}g(zq^k)$ introduces
$\sum_{k\ge0}q^{-ks}=(1-e^{rs})^{-1}$ and gives \eqref{eq:mellin-H}.
\end{proof}

\begin{theorem}[Complete logarithmic asymptotic with Mellin phase]
\label{thm:mellin-asymptotic}
For every fixed $q\in(0,1)$ and every $N>0$, as $z\to+\infty$,
\begin{equation}
 \boxed{
 \log H_q(z)=
 -\frac{L^2}{2r}-\frac L2-\frac r{12}-\frac{c_1}{r}
 +\cP_r(L)+O_{q,N}(z^{-N}).}
 \label{eq:H-asymptotic}
\end{equation}
The periodic function $\cP_r$ is nonconstant.  Therefore no asymptotic that retains
only powers and iterated logarithms can have a relative error tending to zero for
$H_q$ at fixed $q$.
\end{theorem}

\begin{proof}
In \eqref{eq:mellin-H}, shift the contour to the right.  The denominator has simple
poles at $s=\chi_n$; at $n=0$ it combines with the double pole of $g^*$.  The expansion
\[
 \Gamma(1+s)\zeta(1+s)=\frac1s+c_1s+O(s^2),
 \qquad
 \frac1{1-e^{rs}}=-\frac1{rs}+\frac12-\frac{rs}{12}+O(s^3)
\]
gives
\[
 \frac{g^*(s)}{1-e^{rs}}
 =\frac1{rs^3}-\frac1{2s^2}
  +\left(\frac r{12}+\frac{c_1}{r}\right)\frac1s+O(1).
\]
Because the large-$z$ contour displacement contributes the negatives of the crossed
residues, the pole at zero yields the first four nonperiodic terms in
\eqref{eq:H-asymptotic}.  The poles $\chi_n$, $n\ne0$, yield
\eqref{eq:periodic-P}.  The contour may be shifted past any fixed positive real part,
producing $O(z^{-N})$.  Finally, $\Gamma$ has no zeros and the classical zero-free line
$\zeta(1+it)\ne0$ shows that every Fourier coefficient in
\eqref{eq:periodic-P} is nonzero; hence $\cP_r$ is nonconstant.
\end{proof}

\begin{remark}[Why the phase is almost invisible at $q=1/2$]
The first harmonic is suppressed by the gamma factor at
$|\Im\chi_1|=2\pi/r$.  Its exponential scale is $e^{-\pi^2/r}$.  For
$r=\log2$ this is only a few parts in $10^7$ before algebraic factors, whereas for
small $q$ (large $r$) the ripple becomes readily visible.  The attached experiment
plots the residual after removal of the quadratic-logarithmic main part.
\end{remark}

\IfFileExists{figures/mellin-periodic-residual.pdf}{
\begin{figure}[ht]
\centering
\includegraphics[width=0.80\textwidth]{figures/mellin-periodic-residual.pdf}
\caption{The residual in \eqref{eq:H-asymptotic} as a function of $\log z$.  A value of
$q$ for which the Mellin ripple is visible is used in the figure; the dyadic ripple has
the same structure but much smaller amplitude.}
\label{fig:mellin-ripple}
\end{figure}}{}

\section{Endpoint saddle points and phase-aware Lambert--$W$ inversion}
\label{sec:saddle}

The Mellin theorem concerns a transform.  To return to $F_q(x)$ for $x\downarrow0$,
write the inverse Laplace integral in the scaled variable $z=as$:
\begin{equation}
 F_q(x)=\frac1{2\pi i}\int_{c-i\infty}^{c+i\infty}
 e^{(x/a)z}H_q(z)\frac{\Differential z}{z}.
 \label{eq:F-bromwich}
\end{equation}
The saddle itself drifts through the Mellin phase.  Define
\begin{align}
 \cA(L)&=\frac Lr+\frac12-\cP_r'(L),
 \label{eq:A-def}\\
 \cB(L)&=\cA(L)-\cA'(L)
 =\frac Lr+\frac12-\frac1r-\cP_r'(L)+\cP_r''(L).
 \label{eq:B-def}
\end{align}
Both are positive for sufficiently large $L$.

\begin{lemma}[Saddle admissibility condition]
\label{lem:admissibility}
Let $z=e^L$ and exponentially tilt the independent summands in
$\sum q^kZ_k$ by $e^{-z\sum q^kZ_k}$.  Uniformly as $L\to\infty$:
\begin{enumerate}[label=(\alph*)]
\item the tilted variance is $\cB(L)z^{-2}(1+o(1))$;
\item on the central arc $|t|\le z\cB(L)^{-2/5}$, the logarithm of the tilted
characteristic function has a cumulant expansion with standardized cumulants
$O_j(\cB(L)^{1-j/2})$;
\item outside that arc, its modulus contributes $o(\cB(L)^{-1/2})$ to the inversion
integral.
\end{enumerate}
\end{lemma}

\begin{status}
Parts (a) and (b) follow directly by differentiating the product and splitting at
$k\approx L/r$.  Part (c) is the only substantial analytic estimate not written out
line by line here; it follows from the same two-range sinc/product bounds used in
\cref{thm:smoothness}, now for the exponentially tilted factors.  All endpoint saddle
statements below are rigorous conditional on this quantified minor-arc estimate.  The
included code tests the resulting formulas independently at high precision.
\end{status}

\begin{theorem}[Phase-aware endpoint saddle]
\label{thm:phase-saddle}
Assume \cref{lem:admissibility}.  Parameterize $x\downarrow0$ by $L\to\infty$ through
\begin{equation}
 x=x_q(L):=a\cA(L)e^{-L}.
 \label{eq:x-parametric}
\end{equation}
Then
\begin{equation}
 \boxed{
 F_q(x_q(L))=
 \frac{\exp\!\left(
 \cA(L)-\frac{L^2}{2r}-\frac L2-\frac r{12}-\frac{c_1}{r}
 +\cP_r(L)\right)}{\sqrt{2\pi\cB(L)}}
 \left(1+O_q(L^{-1})\right).}
 \label{eq:F-phase-saddle}
\end{equation}
The density satisfies
\begin{equation}
 \frac{f_q(x_q(L))}{F_q(x_q(L))}
 =\frac{e^L}{a}\left(1+O_q(L^{-1})\right).
 \label{eq:hazard-saddle}
\end{equation}
\end{theorem}

\begin{proof}
Let $\ell(z)=\log H_q(z)$.  From \eqref{eq:H-asymptotic},
\[
 -z\ell'(z)=\cA(\log z)+O_N(z^{-N}),
 \qquad
 z^2\ell''(z)=\cB(\log z)+O_N(z^{-N}).
\]
Thus \eqref{eq:x-parametric} is the saddle equation
$x/a=-\ell'(z)$ to beyond every algebraic order.  At the saddle,
$(x/a)z=\cA(L)$ and the exponent in \eqref{eq:F-bromwich} is the exponent displayed in
\eqref{eq:F-phase-saddle}.  Gaussian integration on the central arc gives
$[2\pi\cB(L)]^{-1/2}$; \cref{lem:admissibility} controls the complement and gives the
relative error.  Removing the factor $1/z$ from the inversion integral for the density
multiplies the leading term by $z/a=e^L/a$, proving \eqref{eq:hazard-saddle}.
\end{proof}

\subsection{The Lambert-$W_{-1}$ skeleton}

Suppressing $\cP_r$ in the saddle equation gives
\[
 \frac xa e^L=\frac Lr+\frac12.
\]
With $w=L+r/2$, this is solved exactly by
\begin{equation}
 w=-W_{-1}(-\rho_qx),
 \qquad
 \rho_q=\frac r a e^{-r/2}.
 \label{eq:w-lambert}
\end{equation}
The corresponding phase-free approximation is
\begin{equation}
 F_q(x)\approx
 \sqrt{\frac r{2\pi(w-1)}}
 \exp\!\left[-\frac{w^2-2w}{2r}+\frac r{24}-\frac{c_1}{r}\right].
 \label{eq:F-lambert-skeleton}
\end{equation}
For $q=1/2$, $\rho_q=\sqrt2\log2$.

\begin{remark}[Direct versus inverse asymptotics]
At fixed $q$, omitting $\cP_r$ from \eqref{eq:F-lambert-skeleton} leaves a bounded
nonvanishing multiplicative ripple, so the phase-free direct formula is not an
asymptotic equivalence in the strict ratio sense.  For the inverse function the same
bounded logarithmic ripple shifts $L$ by only $O(1/L)$ because
$\Differential\log F_q/\Differential\log x\sim\cA(L)\asymp L$.  Thus a phase-free inverse can still have
relative error tending to zero, but the phase enters the first genuinely high-accuracy
correction.
\end{remark}

\subsection{A stable parametric inverse}

Define
\begin{equation}
 \mathcal Y_q(L)=
 \cA(L)-\frac{L^2}{2r}-\frac L2-\frac r{12}-\frac{c_1}{r}
 +\cP_r(L)-\frac12\log(2\pi\cB(L)).
 \label{eq:Y-parametric}
\end{equation}
For large $L$, this is strictly decreasing.

\begin{algorithm}[Phase-aware inverse Fabius approximation]
\label{alg:inverse}
Given a sufficiently small $y>0$:
\begin{enumerate}[label=\arabic*.]
\item Solve the scalar equation $\mathcal Y_q(L)=\log y$ for its large positive root.
A phase-free Lambert estimate from \eqref{eq:w-lambert} is an excellent initial value.
\item Return
\begin{equation}
 F_q^{-1}(y)\approx a\cA(L)e^{-L}.
 \label{eq:inverse-parametric}
\end{equation}
\item For higher orders, expand the centered Bromwich integral in Gaussian moments as
described below, add the resulting $L^{-j}$ corrections to
\eqref{eq:Y-parametric}, and repeat one Newton step.
\end{enumerate}
Under \cref{lem:admissibility}, the leading algorithm has relative error
$1+O_q(L^{-2})$ in $F_q^{-1}(y)$.
\end{algorithm}

The improved inverse accuracy follows because an $O(L^{-1})$ error in $\log F_q$ is
divided by the logarithmic slope $\cA(L)=\Theta(L)$.

\subsection{Explicit phase-free inverse expansion}

Let $Y=-\log y$, $S=\sqrt{2rY}$, and
\begin{equation}
 D_r=1+\frac{r^2}{12}-2c_1-r\log\frac{2\pi}{r},
 \qquad A_S=D_r-r\log S.
 \label{eq:Dr-AS}
\end{equation}
Solving the phase-free version of \eqref{eq:Y-parametric} gives
\begin{equation}
 v:=w-1
 =S+\frac{A_S}{2S}
 -\frac{A_S^2+2rA_S}{8S^3}
 +O_r\!\left(\frac{(\log S)^3}{S^5}\right).
 \label{eq:v-expansion}
\end{equation}
Hence
\begin{equation}
 F_q^{-1}(y)=
 \frac{ae^{r/2}}r(v+1)e^{-v-1}
 \left(1+O_r(S^{-1})\right),
 \label{eq:inverse-explicit}
\end{equation}
where the $O(S^{-1})$ term contains the bounded Mellin phase.  Replacing
\eqref{eq:v-expansion} by the phase-aware scalar solve in \cref{alg:inverse} recovers
that term rather than averaging it away.

\subsection{All-orders Bell-polynomial saddle calculus}

The calculation is algorithmic.  Put $z_0=e^L$ and scale the contour by
$z=z_0(1+i\tau/\sqrt{\cB(L)})$.  Then
\[
 \Phi(z)-\Phi(z_0)=-\frac{\tau^2}{2}
 +\sum_{j\ge3}\frac{\Lambda_j(L)}{j!}(i\tau)^j,
 \qquad
 \Lambda_j(L)=\frac{z_0^j\Phi^{(j)}(z_0)}{\cB(L)^{j/2}}.
\]
The amplitude $z_0/z$ is expanded geometrically.  Exponentiating the higher terms uses
partial Bell polynomials; integrating replaces $\tau^{2m}$ by the Gaussian moment
$(2m-1)!!$ and kills odd powers.  Since
\[
 z^j\frac{\Differential^j}{\Differential z^j}
 =D(D-1)\cdots(D-j+1),\qquad D=\frac\Differential{\Differential L},
\]
every coefficient is an explicit differential polynomial in $\cP_r$ and a rational
function of $L$.  The accompanying Python implementation generates these corrections
symbolically or numerically to arbitrary prescribed order.

\section{Consequences for inverse Fabius computations}
\label{sec:inverse-consequences}

For the classical case $q=1/2$, the Mellin phase is tiny but conceptually important.
It changes the interpretation of any purported ``complete'' endpoint expansion.

\begin{proposition}[No phase-free complete direct expansion]
There is no function $A(x)$ built from a finite or asymptotic hierarchy of powers of
$x$, $\log(1/x)$, and iterated logarithms, with no nonconstant periodic dependence on
$\log x/\log2$, such that
\[
 F(x)/A(x)\longrightarrow1\qquad(x\downarrow0)
\]
and such that $A$ captures the saddle to bounded logarithmic error at all orders.
\end{proposition}

\begin{proof}
The transform asymptotic contains the nonconstant periodic function
$\cP_{\log2}(\log z)$.  The saddle map has derivative bounded away from zero after
normalization and transports this term into a nonconstant bounded term in
$\log F(x)$.  A phase-free log-power hierarchy has no term with that periodicity.
The nonvanishing of the Fourier coefficients follows from the proof of
\cref{thm:mellin-asymptotic}.
\end{proof}

\begin{remark}
This does not invalidate useful Lambert-$W$ formulas.  It identifies their exact
asymptotic status.  For ordinary numerical precision, the dyadic phase is usually
smaller than truncation, rounding, and model errors.  For a mathematically complete
transseries or hundreds/thousands of digits, it is unavoidable.
\end{remark}

The phase-aware parametrization also improves numerical conditioning.  Directly
solving $F(x)=y$ near zero involves a function whose logarithmic derivative is large.
In \cref{alg:inverse}, the large variable $L$ is solved from a nearly quadratic,
monotone equation; the final $x=a\cA(L)e^{-L}$ is then formed without subtracting
nearly equal quantities.  The periodic series converges exponentially and usually
requires only one harmonic in the dyadic case.

\section{Legendre coefficients and approximation-theoretic implications}
\label{sec:legendre}

Let $h_b$ be the density in \eqref{eq:Sb}, affinely rescaled to $[-1,1]$ when
necessary, and write its Legendre expansion
\begin{equation}
 h_b(x)=\sum_{n\ge0}a_n^{(b)}P_n(x),
 \qquad
 a_n^{(b)}=\frac{2n+1}{2}\int_{-1}^1h_b(x)P_n(x)\Differential x.
 \label{eq:legendre-expansion}
\end{equation}
Symmetry forces $a_{2n+1}^{(b)}=0$.  Standard repeated-integration estimates for
Legendre coefficients, combined with
\cref{thm:sharp-derivative-growth}, allow the number of integrations to grow with the
degree.

\begin{proposition}[Log-squared upper envelope]
For each $b\ge2$ there is $C_b>0$ such that, along the even degrees,
\begin{equation}
 \log|a_{2n}^{(b)}|
 \le-\frac{(\log n)^2}{2\log b}+C_b\log n\log\log(n+2)
 \qquad(n\ge2).
 \label{eq:legendre-upper}
\end{equation}
The same scale applies to best polynomial approximation errors in every fixed Sobolev
norm.
\end{proposition}

\begin{proof}[Proof sketch]
A $k$-fold Legendre integration-by-parts estimate has the form
$|a_n|\le C^k n^{-k+O(1)}\max_{j\le k}\|h_b^{(j)}\|$; the exact polynomial factors in
$n$ do not affect the leading log-squared term.  By
\cref{thm:sharp-derivative-growth}, the derivative contribution is
$\exp((\log b)k^2/2+O_b(k))$.  Choosing
$k=\Floor{\log n/\log b}$ yields \eqref{eq:legendre-upper}; tracking the
variable-$k$ constants in the Legendre estimate gives the displayed secondary term.
\end{proof}

\begin{conjecture}[Sharp Legendre envelope]
\label{conj:legendre-envelope}
For every integer $b\ge2$,
\begin{equation}
 \limsup_{n\to\infty}
 \frac{\log|a_{2n}^{(b)}|}{(\log n)^2}
 =-\frac1{2\log b}.
 \label{eq:legendre-envelope-conj}
\end{equation}
More precisely, after removal of the quadratic and linear terms in $\log n$, the
upper envelope contains a bounded nonconstant periodic function of
$\log n/\log b$.
\end{conjecture}

The conjectured constant is the Legendre-dual of the derivative-growth constant:
optimizing $n^{-k}b^{k^2/2}$ produces exactly
$\exp[-(\log n)^2/(2\log b)]$.  A lower bound is harder because Legendre oscillations
can create cancellations absent from the Fourier resonant sequence.  The supplied code
computes exact rational moments for rational $q$, forms Legendre coefficients without
quadrature, and plots the renormalized envelope.

\IfFileExists{figures/legendre-envelope.pdf}{
\begin{figure}[ht]
\centering
\includegraphics[width=0.80\textwidth]{figures/legendre-envelope.pdf}
\caption{Renormalized even Legendre coefficients for the dyadic density.  The plot is a
diagnostic for \cref{conj:legendre-envelope}; finite degrees should not be mistaken for
proof of the limiting constant.}
\label{fig:legendre-envelope}
\end{figure}}{}

\section{Further conjectures and research programs}
\label{sec:future}

\subsection{Exponentially improved endpoint transseries}

\begin{conjecture}[Geometric exponential transseries]
After the full periodic saddle series in powers of $1/L$ is removed, the endpoint
remainder admits a transseries generated by exponentials
\[
 e^{-zq^j},\qquad j=0,1,2,\dots,
\]
where $z=e^L$ is the saddle variable.  The sectors accumulate geometrically and obey a
$q$-difference resurgence law inherited from
$H_q(z)=((1-e^{-z})/z)H_q(qz)$.
\end{conjecture}

The superalgebraic remainder in \cref{thm:mellin-asymptotic} is consistent with this
claim but does not identify its exponentially small structure.  A useful first goal is
to derive the leading $e^{-z}$ correction uniformly through one Mellin period and then
track its descendants under $z\mapsto qz$.

\subsection{Minor-arc closure of the saddle theorem}

\begin{problem}
Give a self-contained proof of \cref{lem:admissibility}(c) with an explicit constant
uniform for $q$ in compact subsets of $(0,1)$.  Then compute the first three relative
saddle coefficients and certify them by interval arithmetic.
\end{problem}

A direct route is to split the tilted product into active indices
$k\le L/r+O(1)$ and a small-argument tail.  On the active block, a positive proportion
of factors stays uniformly away from its local lattice of maxima; on the tail, the
quadratic characteristic-function estimate is uniform.  This is analogous to the
minor-arc component of local limit theorems for digital harmonic sums.

\subsection{Zero divisors beyond reciprocal integers}

For general $q$, the multiplicity function is
\[
 m_q(z)=\#\{k\ge0:zq^k/\pi\in\IntegerNumbers\setminus\{0\}\}.
\]
When $q^d=A/B$ is rational in lowest terms, the layers split into $d$ residue classes,
each carrying a mixed $A$--$B$ valuation structure.

\begin{problem}
Classify the canonical product and spectral Dirichlet series for all $q$ such that some
$q^d$ is rational.  Determine whether the resulting Dirichlet series factor into
Riemann zeta functions and finitely many local Euler corrections, or whether genuinely
new automatic Dirichlet series occur when the numerator $A>1$.
\end{problem}

\subsection{Geometric Bernoulli zeros and operators}

The averaging operator
\[
 (T_qp)(x)=\Expectation p(x+V_q)
\]
has inverse on polynomials given by $T_q^{-1}x^n=\Bell_{n,q}(x)$.  The fixed-point
decomposition factors it as an infinite product of uniform-averaging operators.

\begin{problem}
Develop the spectral and zero geometry of $\Bell_{n,q}$.  Questions include the
location of nonreal zeros, interlacing under $n\mapsto n+1$, asymptotic zero measures as
$n\to\infty$, and operator factorizations of $T_q^{-1}$ in terms of the $q$-difference
law \eqref{eq:q-difference-Bernoulli}.
\end{problem}

\begin{conjecture}[Central zero rigidity]
For each fixed $q\in(0,1)$, all real zeros of $\Bell_{n,q}$ are simple except for the
forced central zero of odd degree, and the number of real zeros is eventually a fixed
proportion of $n$ depending continuously on $q$.
\end{conjecture}

This is deliberately stated as a numerical research question rather than a theorem;
classical Bernoulli polynomial zeros already warn against naive all-real-rootedness.

\subsection{Strict log-concavity and total positivity}

\begin{conjecture}[Strict interior log-concavity]
For every $q\in(0,1)$, $(\log u_q)''(x)<0$ for $-1<x<1$, except possibly at the
symmetry point in a limiting interpretation.
\end{conjecture}

Finite convolutions are spline densities for which strictness can be analyzed by
variation-diminishing methods.  Passing a uniform strict estimate to the infinite
limit would give strong inequalities for $F_q$, including monotonicity of endpoint
hazard ratios used in rigorous inverse bounds.

\subsection{Formalization targets}

Several results are well suited for formal proof assistants:
\begin{enumerate}[label=(\roman*)]
\item the probabilistic fixed point and functional equations;
\item the affine-word induction in \cref{thm:qthue-derivative};
\item the Bell-polynomial formula \eqref{eq:Snm-bell};
\item the moment recurrence and geometric Bernoulli inversion;
\item the valuation Hadamard product at the level of formal power series; and
\item the upper half of the derivative-growth theorem.
\end{enumerate}
The resonant lower bound additionally needs a modest amount of real/complex analysis.
The Mellin contour argument is the least immediate formalization target but can be
split into reusable theorems on harmonic sums.

\section{Numerical experiments and reproducibility}
\label{sec:numerics}

The archive contains \texttt{experiments.py}.  It uses arbitrary-precision arithmetic
for products and Mellin coefficients, exact rational arithmetic where feasible, and
standard floating-point quadrature only for the optional density plots.  Running
\begin{center}
\texttt{python experiments.py --all}
\end{center}
regenerates the CSV tables and PDF figures used here.

The experiments are designed as tests of identities, not as substitutes for proof:
\begin{enumerate}[label=(\arabic*)]
\item \texttt{mellin\_comparison.csv} compares the direct product
$\log H_q(z)$ with \eqref{eq:H-asymptotic}, both with and without the periodic term.
\item \texttt{resonant\_fourier.csv} checks \eqref{eq:resonant-fourier} at high
precision for several bases and levels.
\item \texttt{edgeworth\_comparison.csv} compares numerical Fourier inversion with
\eqref{eq:edgeworth-density}.
\item \texttt{legendre\_coefficients.csv} computes dyadic even Legendre coefficients
from exact moments, avoiding endpoint quadrature.
\item \texttt{corpus-manifest.txt} records every audited TeX source and its hash.
\end{enumerate}

\IfFileExists{data/mellin_comparison.tex}{\input{data/mellin_comparison.tex}}{}

\section{Synthesis}

The geometric uniform convolution turns several themes that appear separately in the
classical theory into facets of a single object.
\begin{itemize}
\item The two-branch differential refinement generates Thue--Morse signs; the branch
locations carry a geometric $q$-deformation, and their signed moments are controlled by
Bell polynomials with Bernoulli-number inputs.
\item The infinite version of the same geometric power sums gives the exact cumulants.
The reciprocal moment-generating function then produces an Appell family that extends
classical Bernoulli polynomials and exactly inverts Fabius averaging.
\item At reciprocal integer $q$, the Fourier zeros merge according to a digital
valuation.  The same base $b$ fixes both the zero multiplicities and the sharp
$n^2\log b/2$ derivative-growth scale.
\item Mellin poles at an imaginary arithmetic progression record discrete scale
invariance.  Their log-periodic term survives the saddle and supplies the missing phase
in a genuinely complete endpoint expansion.
\item Lambert $W_{-1}$ remains the correct coarse inverse coordinate, but the Mellin
phase specifies how to pass from a highly accurate approximation to an asymptotic
calculus that is complete at fixed $q$.
\end{itemize}

The strongest concrete result is \cref{thm:sharp-derivative-growth}: for the classical
up-function, the exact leading logarithmic growth of every $L^p$ derivative norm is
$(\log2)n^2/2$.  The most useful new computational object is the phase-aware inverse
parametrization \eqref{eq:Y-parametric}--\eqref{eq:inverse-parametric}.  The broadest
new structural object is the geometric Bernoulli sequence, which may be developed
independently as a $q$-Appell/N\o rlund theory.

\appendix

\section{Selected algebraic checks}

\subsection{Consistency of the zeta and Bernoulli expansions}

Euler's formula
\[
 \zeta(2m)=(-1)^{m+1}\frac{B_{2m}(2\pi)^{2m}}{2(2m)!}
\]
turns \eqref{eq:Psi-Taylor-zeta} into the cumulant expansion obtained from
\eqref{eq:X-cumulants} after the rescaling $z=at$.  This independently checks all
normalizing factors in the Fourier and probabilistic derivations.

\subsection{The first nonvanishing geometric Prouhet moment}

Near $t=0$,
\[
 2^n\prod_{j=0}^{n-1}\sinh(tq^j)
 =2^nt^nq^{n(n-1)/2}+O(t^{n+2}),
\]
so differentiation at zero gives
$S_{n,n}=2^nn!q^{n(n-1)/2}$.  Substituting this into the Taylor expansion of
\eqref{eq:qthue-derivative} confirms that the normalization reproduces the $n$th
derivative rather than a scalar multiple.

\subsection{Scaling invariance of the quadratic derivative constant}

If $g(x)=c^{-1}f(x/c)$, then
\[
 \|g^{(n)}\|_p=c^{-n-1+1/p}\|f^{(n)}\|_p.
\]
The logarithm of the scaling factor is linear in $n$ and therefore disappears after
division by $n^2$.  This justifies moving freely between $h_b$ and the support-normalized
$u_{1/b}$ in \cref{thm:sharp-derivative-growth}.

\section{Corpus manifest}

The full manifest is supplied as a separate text file in the archive to keep this PDF
readable.  A compact generated summary is included below when available.

\IfFileExists{corpus-summary.tex}{\input{corpus-summary.tex}}{
\begin{quote}
Corpus summary was not generated in this build.  See \texttt{corpus-manifest.txt}.
\end{quote}}

\begin{thebibliography}{99}

\bibitem{Fabius1966}
J.~Fabius,
\newblock A probabilistic example of a nowhere analytic $C^\infty$-function,
\newblock \emph{Zeitschrift f\"ur Wahrscheinlichkeitstheorie und Verwandte Gebiete}
\textbf{5} (1966), 173--174.

\bibitem{Rvachev}
V.~L. Rvachev,
\newblock Compactly supported solutions of functional-differential equations and their
applications,
\newblock foundational papers and monographs on atomic functions; bibliographic details
and translations are collected in the audited repository corpus.

\bibitem{deBruijn1948}
N.~G. de Bruijn,
\newblock On Mahler's partition problem,
\newblock \emph{Indagationes Mathematicae} \textbf{10} (1948), 210--220.

\bibitem{FlajoletGourdonDumas1995}
P.~Flajolet, X.~Gourdon, and P.~Dumas,
\newblock Mellin transforms and asymptotics: harmonic sums,
\newblock \emph{Theoretical Computer Science} \textbf{144} (1995), 3--58.

\bibitem{AlloucheShallit}
J.-P. Allouche and J.~Shallit,
\newblock \emph{Automatic Sequences: Theory, Applications, Generalizations},
\newblock Cambridge University Press, 2003.

\bibitem{GasperRahman}
G.~Gasper and M.~Rahman,
\newblock \emph{Basic Hypergeometric Series}, 2nd ed.,
\newblock Cambridge University Press, 2004.

\bibitem{Comtet}
L.~Comtet,
\newblock \emph{Advanced Combinatorics},
\newblock D.~Reidel, 1974.

\bibitem{Norlund}
N.~E. N\o rlund,
\newblock \emph{Vorlesungen \"uber Differenzenrechnung},
\newblock Springer, 1924.

\bibitem{Petrov}
V.~V. Petrov,
\newblock \emph{Sums of Independent Random Variables},
\newblock Springer, 1975.

\bibitem{Wong}
R.~Wong,
\newblock \emph{Asymptotic Approximations of Integrals},
\newblock SIAM, 2001.

\bibitem{Rodino}
L.~Rodino,
\newblock \emph{Linear Partial Differential Operators in Gevrey Spaces},
\newblock World Scientific, 1993.

\bibitem{Szego}
G.~Szeg\H{o},
\newblock \emph{Orthogonal Polynomials}, 4th ed.,
\newblock American Mathematical Society Colloquium Publications, vol.~23, 1975.

\bibitem{DLMF}
F.~W.~J. Olver, D.~W. Lozier, R.~F. Boisvert, and C.~W. Clark (eds.),
\newblock \emph{NIST Handbook of Mathematical Functions},
\newblock Cambridge University Press, 2010.

\bibitem{ProveItCorpus}
V.~Reshetnikov,
\newblock \emph{ProveIt: Analysis/FabiusFunction/docs},
\newblock Git repository, snapshot \RepoCommit, recursively audited for this report.

\end{thebibliography}

\end{document}
